%-------------------------------------------------------------

\medskip

\begin{enumerate}
	\item D'une part $AD^2=7,3^2=53,29$.\\
	D'autre part $AE^2+ED^2=5,5^2+4,8^2=53,29$\\
	Donc $AD^2=AE^2+ED^2$\\
	D'après la réciproque du théorème de Pythagore, on conclut que le triangle AED est rectangle en E.
	\item $\mathcal{A}_{AED}=\dfrac{5,5\times 4,8}{2}=13,2$ cm\up{2}.
	\item Les droites $(CB)$ et $(ED)$ sont toutes les deux perpendiculaires à la même droite $(BE)$ donc elles sont parallèles entre elles.
	\item On sait que : \begin{itemize}
							\item $(CB)\perp (ED)$ ; 
							\item les points $E,~A,~B$ et $D,~A,~C$ sont alignés dans le même ordre.
						\end{itemize}	  
						D'après le théorème de Thalès, on a :
						\begin{align*}
						\dfrac{AE}{AB}&=\dfrac{AD}{AC}=\dfrac{ED}{BC}\\
						\dfrac{5,5}{AB}&=\left(\dfrac{7,3}{AC}\right)=\dfrac{4,8}{7,2}\\
						AB&=\dfrac{7,2\times 5,5}{4,8}\\
						AB&=\dfrac{33}{4}=8,25 \text{ cm}
						\end{align*}
	\item Les longueurs des côtés des triangles $ABC$ et $AED$ sont proportionnelles donc les triangles sont semblables. Donc leurs angles sont deux à deux de même mesure. Ainsi $\widehat{ADE}=\widehat{ACB}\approx 49^\circ$.
\end{enumerate}

\bigskip

%-------------------------------------------------------------
